
\section{Conclusion and Outlook} \label{conclusion}

In conclusion, in this project a complete end-to-end 3D multi-person tracking pipeline was designed, implemented, and evaluated.
The quantitative evaluation on the CMU Panoptic dataset provided some architectural insight. Spatial viewpoint density exhibits a far higher return on tracking accuracy than increasing deep learning model capacity. Scaling the camera view count from 4 to 16 with the YOLO26n-pose model reduced the mean per-joint position error from 8.25 cm down to 2.83 cm. In contrast, replacing this lightweight model with the heavy YOLO26l-pose model under a fixed 4-camera setup yielded smaller accuracy gains at the cost of a massive computational bottleneck, dropping system throughput to approximately 3 FPS.
Additionally, performance can also be significantly improved with the addition of skeletal constraints as well as the use of learning from data for different components.  

Despite the relatively high localization accuracy of the system under optimal conditions, there remain several limitations. As demonstrated in the qualitative evaluations, tracking quality degrades when targets are partially occluded or not in view of the camera views. Since 3D initialization requires triangulation across at least two views, the system also exhibits tracking delays during entry and spatial drifting during single-camera visibility. Further, the way the system has been adapted to the dataset by training individual components separately limits the amount of adaptation.

Possible future improvements to address these limitations could include fully end-to-end training of the system parameters from the dataset instead of training components separately. A further improvement could be the addition of further anatomical plausibility constraints on the state predictions, such as realistic joint angles. These constraints could possibly be included similar to the existing limb length constraints using a pseudo-measurement. Also, the addition of additional information sources, such as monocular depth values could make it possible to track people even using single camera views. Exploring more advanced filtering methods such as particle filters, potentially also using occlusion aware likelihoods could be promising future directions. Finally, the evaluation in this project has been performed only on the CMU Panoptic dataset. As such, the evaluation of the system using a larger set of more diverse sequences would be highly valuable to uncover further areas for improvement.
